\documentclass[11pt]{article}
\usepackage[T1]{fontenc}
\usepackage{amsmath,amssymb,amsthm}
\usepackage[margin=1in]{geometry}
\usepackage[colorlinks=true,linkcolor=blue]{hyperref}
\newtheorem{theorem}{Theorem}
\newcommand{\R}{\mathbb{R}}
\title{The Spectral Theorem}
\author{MatWiki}
\date{}
\begin{document}
\maketitle

\begin{theorem}
Every real symmetric matrix $A\in\R^{n\times n}$ has an orthonormal basis
of real eigenvectors. Equivalently, $A=Q\Lambda Q^{\mathsf T}$ for an
orthogonal matrix $Q$ and a real diagonal matrix $\Lambda$.
\end{theorem}

\begin{proof}
Induct on $n$. The result is immediate for $n=1$. The real function
$x\mapsto x^{\mathsf T}Ax$ attains a maximum on the unit sphere.
By Lagrange multipliers, a maximiser $v$ satisfies $Av=\lambda v$.
For $w\perp v$, symmetry gives
\[
  \langle Aw,v\rangle=\langle w,Av\rangle
  =\lambda\langle w,v\rangle=0.
\]
Thus $v^\perp$ is invariant under $A$, and the restriction is symmetric.
Apply the induction hypothesis on $v^\perp$ and append $v$ to the resulting
orthonormal eigenbasis.
\end{proof}

\section*{Example}
\[
  A=\begin{pmatrix}2&1\\1&2\end{pmatrix},\qquad
  Q=\frac{1}{\sqrt{2}}\begin{pmatrix}1&1\\1&-1\end{pmatrix},\qquad
  \Lambda=\begin{pmatrix}3&0\\0&1\end{pmatrix}.
\]

% Wiki connections live in comments so the document compiles cleanly.
% [[linear-algebra]]
% [[00-start-here]]
\end{document}
